\section{Implementierung}

Dieses Kapitel beschreibt die technische Umsetzung des Bestellanforderungsprozesses. Die Implementierung erfolgt auf Basis des \gls{cap} mit Node.js, SAP Fiori Elements und \gls{spa}. Vollständige Code-Listings befinden sich im Anhang.

\subsection{Projektstruktur}

Die Anwendung folgt der \gls{cap}-Konvention mit drei Hauptverzeichnissen: \texttt{db/} enthält das \gls{cds}-Domänenmodell, \texttt{srv/} die Servicedefinition und Geschäftslogik, \texttt{app/} die Fiori-Elements-Anwendung. Die Dateien \texttt{mta.yaml} und \texttt{xs-security.json} steuern Deployment und Sicherheit.

\subsection{Datenbankschicht}

Das Domänenmodell definiert zwei Hauptentitäten in \gls{cds} (siehe Anhang~A). Die Entität \texttt{PurchaseRequisitions} repräsentiert Bestellanforderungen mit Beschreibung, Anforderer, Kostenstelle, Gesamtbetrag und Status (\texttt{Draft}, \texttt{Submitted}, \texttt{Approved}, \texttt{Rejected}). Die Entität \texttt{Items} enthält die zugehörigen Positionen mit Materialnummer, Menge, Einzelpreis und berechnetem Nettobetrag. Die Beziehung ist als Komposition modelliert, sodass Positionen gemeinsam mit der übergeordneten Anforderung verwaltet werden. Zusätzlich protokolliert eine \texttt{StatusHistory}-Entität alle Statusänderungen, und das Plugin \texttt{@cap-js/attachments} ermöglicht Dokumentenanhänge.

\subsection{Service-Schicht}

Die Servicedefinition (siehe Anhang~B) exponiert das Domänenmodell als \gls{odata}-V4-Service mit drei Kernfunktionen: Die Draft-Unterstützung durch die Annotation \texttt{@odata.draft.enabled} ermöglicht das schrittweise Bearbeiten von Anforderungen. Die Submit-Aktion validiert die Anforderung und startet den Genehmigungsworkflow. Der Callback-Endpunkt empfängt die Genehmigungsentscheidung von \gls{spa}. Die Geschäftslogik in den Event-Handlern berechnet automatisch Netto- und Gesamtbeträge beim Speichern und protokolliert alle Statusänderungen in der Historie.

\subsection{Workflow-Integration}

Die Kommunikation mit \gls{spa} erfolgt über einen dedizierten \texttt{WorkflowClient}, der den OAuth2 Client Credentials Flow implementiert. Beim Freigeben einer Anforderung wird ein Genehmigungsworkflow gestartet, der folgende Kontextdaten erhält: Anforderungs-ID, Beschreibung, Anforderer, Gesamtbetrag, Positionsliste, Anhänge-URLs sowie ein Deep Link zur Fiori-Anwendung.

Nach der Genehmigungsentscheidung ruft \gls{spa} den Callback-Endpunkt auf und übermittelt Status, Genehmigername und Kommentar. Die CAP-Anwendung aktualisiert daraufhin den Status der Bestellanforderung.

\subsection{Frontend}

Die Benutzeroberfläche basiert vollständig auf SAP Fiori Elements und wird deklarativ über \gls{cds}-Annotationen konfiguriert. Der List Report zeigt alle Bestellanforderungen tabellarisch an, die Object Page ermöglicht die Detailbearbeitung mit Sektionen für Kopfdaten, Positionen, Genehmigungsinformationen, Statusverlauf und Anhänge.

Eine benutzerdefinierte Schaltfläche ermöglicht die Navigation zur SPA My Inbox, wo Genehmiger ihre Aufgaben bearbeiten können.

\subsection{Deployment}

Die Bereitstellung erfolgt als \gls{mta} auf Cloud Foundry (siehe Anhang~D). Das Deployment umfasst drei Module: Der CAP-Server ist eine Node.js-Anwendung mit \gls{odata}-Service und Geschäftslogik. Der HDI-Deployer überführt das CDS-Modell in das HANA-Datenbankschema. Der Application Router liefert die Fiori-Oberfläche aus und koordiniert die Authentifizierung. Die \gls{spa}-Konfiguration erfolgt über Umgebungsvariablen, was die Trennung von Code und Konfiguration nach dem Twelve-Factor-App-Prinzip gewährleistet \citep{wiggins2017twelvefactor}.
